\capitulosinnumero{Justificación y objetivos}

El presente Trabajo Fin de Grado surge de la necesidad de explorar nuevas formas de expresión interactiva en el medio del videojuego que integren el discurso narrativo directamente en sus mecánicas jugables. Aunque la industria independiente ha avanzado en la representación de temáticas introspectivas y de salud mental, gran parte de las obras actuales continúan separando la historia de la jugabilidad o implementan sistemas dinámicos que alteran las reglas de forma unilateral, generando frustración y pérdida de agencia en el usuario.
Desde una perspectiva social, académica y técnica propia del Grado en Ingeniería Multimedia, este proyecto busca visibilizar el agotamiento psicológico derivado de una crisis de identidad de género a causa del rechazo social e internalizado. Para ello, se plantea el diseño y desarrollo en el motor Unity de un videojuego narrativo en 2D con vista cenital oblicua (\textit{top-down} 3/4) que emplea un sistema de mecánicas adaptativas en segundo plano y una mecánica de consentimiento explícito.
El objetivo general del proyecto consiste en desarrollar un videojuego narrativo en 2D con un sistema de mecánicas adaptativas que explore una crisis de identidad desde la perspectiva de género, desglosándose en la elaboración del Documento de Diseño de Juego (\ac{GDD}) y su dirección artística con asesoramiento psicológico, la programación del módulo evaluador del comportamiento en segundo plano, la creación de una arquitectura flexible basada en \textit{ScriptableObjects} para mutar habilidades e interfaz (\ac{UI}) en tiempo real, y la construcción de un prototipo funcional y jugable.